\section{Analysis and Limitations}

Three design lessons emerge from our evaluation. Tailoring input text to specific scoring mechanisms allows the hybrid interface to leverage exact catalog cues while the conversational interface captures evolving language constraints, as illustrated in Figure~\ref{fig:textviews}. Furthermore, the low candidate overlap and complementary target hits between these distinct channels create productive pool disagreement, explaining why cross-pool RRF outperforms single-pool baselines in Section~\ref{sec:experiments}. Finally, structuring response generation through explicit evidence and role assignments ensures that PAS prioritizes evidence-backed track selection over surface realization.

These findings outline trajectories for architectural extension, particularly regarding the LightGBM ranker that currently relies on two upstream pools and compact features. Integrating stronger retrieval components with diverse behaviors will provide more informative signals to the calibrator. Incorporating richer features that directly relate each track to the query and dialogue trajectory offers a direct path to maximize reranking gains.

Several limitations remain in the current implementation. Uneven external metadata can constrain PAS despite deterministic validation, and the in-house judge may introduce model bias during offline compilation. Additionally, the 18-example synthetic bank cannot cover every dialogue state or sparse request bucket. The five adapted 8B folds and LLM-based synthesis also require substantial compute. Consequently, this study leaves runtime latency, serving costs, and human preferences unevaluated.